\documentclass[a4paper,12pt]{article}

%% ── Geometry (EPO A4 standard margins) ──
\usepackage[a4paper,top=25mm,bottom=25mm,left=25mm,right=25mm]{geometry}

%% ── Fonts and Encoding ──
\usepackage[T1]{fontenc}
\usepackage[utf8]{inputenc}

%% ── Mathematics ──
\usepackage{amsmath,amssymb}

%% ── Graphics ──
\usepackage{graphicx}
\usepackage{tikz}
\usetikzlibrary{arrows.meta,positioning,shapes.geometric,fit,calc}

%% ── Tables ──
\usepackage{booktabs}
\usepackage{tabularx}
\usepackage{array}

%% ── Hyperlinks ──
\usepackage[colorlinks=false]{hyperref}

%% ── Headers/Footers ──
\usepackage{fancyhdr}
\pagestyle{fancy}
\fancyhf{}
\fancyhead[L]{\small\textit{COGNIGRAPH-CIP-2026 --- Description}}
\fancyhead[R]{\small\textit{Quantamix Solutions B.V.}}
\fancyfoot[C]{\small\thepage}
\renewcommand{\headrulewidth}{0.4pt}
\renewcommand{\footrulewidth}{0.2pt}

%% ── Compact lists ──
\usepackage{enumitem}

%% ── Line spacing ──
\usepackage{setspace}
\setstretch{1.15}

\begin{document}

\begin{center}
{\large\bfseries DESCRIPTION}\\[3mm]
{\normalsize European Patent Application --- Continuation-in-Part}\\[1mm]
{\normalsize Parent: EP26162901.8 (filed 2026-03-06)}\\[1mm]
{\normalsize Reference: COGNIGRAPH-CIP-NL-2026-001}\\[1mm]
{\normalsize Applicant: Quantamix Solutions B.V.}\\[3mm]
{\large\bfseries CogniGraph: Graph-of-Agents Distributed Reasoning\\[2mm]
over Knowledge Graph Topologies with Prize-Collecting\\[2mm]
Steiner Tree Activation and Convergent Message Passing}
\end{center}

\vspace{8mm}

%% =============================================================================
%% SECTION 1: TECHNICAL FIELD
%% =============================================================================

\section*{Technical Field}

\noindent The present invention relates to the field of artificial intelligence and knowledge graph reasoning, and more particularly to a system and method for distributed multi-agent reasoning over knowledge graph topologies using graph-theoretic subgraph activation, convergent message passing, and hierarchical aggregation.

\vspace{4mm}

%% =============================================================================
%% SECTION 2: RELATIONSHIP TO PARENT APPLICATION
%% =============================================================================

\section*{Relationship to Parent Application}

\noindent This application is a Continuation-in-Part (CIP) of European Patent Application EP26162901.8, entitled ``TAMR+: Trust-Aware Multi-Signal Document Retrieval with Graph-Based Compliance Scoring and Gap Attribution for Regulatory AI Systems,'' filed on 6~March~2026 by Quantamix Solutions B.V.

\vspace{2mm}

\noindent The parent application EP26162901.8 discloses a document retrieval and compliance scoring system (TAMR+) that uses Prize-Collecting Steiner Tree (PCST) optimization for knowledge graph traversal in the context of regulatory document retrieval. The present CIP application extends the PCST activation concept from document retrieval to \textbf{distributed multi-agent reasoning}, introducing five novel innovations:

\begin{enumerate}[nosep,leftmargin=2em]
\item \textbf{PCST-Activated Multi-Agent Reasoning} --- using PCST subgraph activation to determine agent deployment topology (parent: PCST for retrieval subgraph selection);
\item \textbf{MasterObserver Transparency Monitoring} --- a non-participant observer that monitors inter-agent reasoning traffic for transparency and auditability;
\item \textbf{Convergent Message Passing} --- similarity-based convergence detection with early termination and token compression;
\item \textbf{Backend Fallback Chain with Cost Budget} --- resilient inference across heterogeneous model backends with per-query cost enforcement;
\item \textbf{Graph-Topology-Aware Hierarchical Aggregation} --- centrality-based aggregation tree construction for synthesizing distributed reasoning results.
\end{enumerate}

\vspace{2mm}

\noindent \textbf{Shared subject matter} (inheriting parent priority date): The PCST optimization algorithm for subgraph selection on knowledge graphs is common to both applications. Claims 1--3 of this CIP extend the PCST activation from the parent's retrieval context to the novel reasoning context.

\noindent \textbf{New matter} (receiving CIP filing date): Innovation Groups B--E (Claims 4--15) constitute new matter not disclosed in the parent application, covering multi-agent monitoring, convergent message passing, backend resilience, and hierarchical aggregation.

\vspace{4mm}

%% =============================================================================
%% SECTION 3: BACKGROUND OF THE INVENTION
%% =============================================================================

\section*{Background of the Invention}

\noindent Knowledge graphs encode structured domain knowledge as nodes (entities) connected by typed edges (relationships). Reasoning over knowledge graphs has traditionally been performed by (a)~graph traversal algorithms that follow edges to find relevant paths, (b)~embedding-based methods that project entities and relations into vector spaces, or (c)~retrieval-augmented generation (RAG) that extracts subgraph context and feeds it to a single language model.

\vspace{2mm}

\noindent These approaches share a fundamental limitation: reasoning is \textbf{centralized}. A single model receives all relevant context and must reason over it in one pass. This creates three problems:

\begin{enumerate}[nosep,leftmargin=2em]
\item \textbf{Context window saturation}: Large knowledge graphs produce context that exceeds model capacity;
\item \textbf{Loss of graph structure}: Concatenating node content destroys the topological relationships that encode domain semantics;
\item \textbf{No transparency}: A single model's reasoning is opaque --- there is no mechanism to trace which knowledge sources contributed to which parts of the answer.
\end{enumerate}

\vspace{2mm}

\noindent The parent application (TAMR+) addresses the retrieval component of this pipeline by using PCST to select relevant subgraphs for document retrieval and compliance scoring. However, TAMR+ feeds retrieved content to a single language model for final reasoning, inheriting the centralization limitations above.

\vspace{2mm}

\noindent CogniGraph addresses the reasoning component by transforming the knowledge graph itself into a \textbf{reasoning topology}: each node becomes an autonomous language model agent that reasons about its local knowledge and communicates with neighbor agents through the graph's edge structure. This Graph-of-Agents architecture preserves graph topology during reasoning, distributes context across multiple agents, and provides full transparency through the MasterObserver.

\vspace{4mm}

%% =============================================================================
%% SECTION 4: SUMMARY OF THE INVENTION
%% =============================================================================

\section*{Summary of the Invention}

\noindent The present invention provides a system and method for distributed reasoning over knowledge graphs, called CogniGraph, comprising:

\begin{enumerate}[nosep,leftmargin=2em]
\item \textbf{Graph-of-Agents Architecture}: Each node $v_i$ in a knowledge graph $G = (V, E, W)$ hosts an autonomous language model agent $a_i$ that has access to the node's knowledge content, entity type, and local graph neighborhood.

\item \textbf{PCST Subgraph Activation}: For each query $q$, a Prize-Collecting Steiner Tree optimization selects a compact subgraph $G^* \subseteq G$ of 4--$K_{\max}$ nodes, achieving sublinear scaling with graph size.

\item \textbf{Convergent Message Passing}: Agents conduct $R$ rounds of message-passing reasoning, where each agent generates responses conditioned on the query, its node knowledge, and messages from graph-adjacent agents. Convergence is detected via inter-message similarity, enabling early termination.

\item \textbf{MasterObserver}: A non-reasoning observer module intercepts all inter-agent traffic, providing complete transparency, health monitoring, and contribution attribution without affecting reasoning outcomes.

\item \textbf{Backend Fallback Chain}: Agents can use heterogeneous inference backends (cloud APIs, local GPU models, adapter-augmented models) with automatic failover and per-query cost budget enforcement.

\item \textbf{Hierarchical Aggregation}: Final-round messages are aggregated using a centrality-based tree structure (leaves $\to$ hubs $\to$ root) that preserves the graph's topological hierarchy in the synthesis.
\end{enumerate}

\vspace{4mm}

%% =============================================================================
%% SECTION 5: DETAILED DESCRIPTION
%% =============================================================================

\section*{Detailed Description of the Invention}

\subsection*{5.1 System Architecture}

\noindent The CogniGraph system comprises five principal components:

\begin{enumerate}[nosep,leftmargin=2em]
\item \textbf{Knowledge Graph Layer}: Accepts graph data from NetworkX, Neo4j, or JSON/node-link format. Each node stores: unique identifier, label, entity type, description (knowledge content), and typed properties. Edges store: source/target identifiers, relationship type, and weight.

\item \textbf{Activation Layer}: Implements PCST subgraph selection. Given query $q$, computes node prizes $\pi_i = \alpha \cdot \mathrm{sim}(\mathbf{e}_q, \mathbf{e}_{v_i})$ using sentence embeddings, and edge costs $c_e = \beta / w_e$. The PCST solver (Goemans--Williamson primal-dual) returns the activated subgraph $G^*$.

\item \textbf{Orchestration Layer}: Manages message-passing rounds, convergence detection, and cost tracking. For each round $r$, agents generate responses in parallel; the orchestrator collects messages, checks convergence ($\sigma^{(r)} \geq \tau$), and decides whether to continue or terminate.

\item \textbf{Observation Layer}: The MasterObserver intercepts all messages via an event bus, recording trace data without modifying the reasoning flow. Computes health scores, contribution metrics, and anomaly alerts.

\item \textbf{Backend Layer}: Provides a unified interface to heterogeneous language model backends. Implements the fallback chain $(b_1, b_2, \ldots, b_k)$ with retry logic, exponential backoff, and cost tracking.
\end{enumerate}

\subsection*{5.2 PCST Activation for Agent Deployment}

\noindent Unlike the parent application where PCST selects document nodes for retrieval, CogniGraph uses PCST to determine the \textbf{reasoning topology} --- which nodes should host active agents and which edges should carry messages. This is a fundamentally different use of the same optimization:

\begin{itemize}[nosep,leftmargin=2em]
\item \textbf{TAMR+ (parent)}: PCST selects nodes whose content should be retrieved and concatenated into a prompt for a single model.
\item \textbf{CogniGraph (this CIP)}: PCST selects nodes that become active reasoning agents, preserving the edge structure as communication channels between agents.
\end{itemize}

\noindent The prize function and cost function use the same mathematical formulation as the parent application, but the output is interpreted differently: instead of a retrieval set, it defines an agent deployment topology.

\subsection*{5.3 Message-Passing Reasoning Protocol}

\noindent Each activated agent $a_i$ maintains a context comprising:

\begin{enumerate}[nosep,leftmargin=2em]
\item A \textbf{system prompt} derived from the node's entity type and knowledge content:
\begin{quote}
\small\ttfamily
You are a reasoning agent for node [\textit{label}] ([\textit{type}]). \\
Your knowledge: [\textit{description}]. \\
Reason about the query using your knowledge and messages from neighbors.
\end{quote}

\item The \textbf{user query} $q$;
\item \textbf{Incoming messages} from graph-adjacent agents: $\{m_j^{(r-1)} : (v_j, v_i) \in E^*\}$.
\end{enumerate}

\noindent In each round $r$, the agent generates response $m_i^{(r)}$ by invoking its backend model with this context. Messages exceeding the token budget $B_{\mathrm{tok}}$ are compressed using the $\mathrm{compress}(\cdot)$ function described in Claim~9.

\subsection*{5.4 Convergence Detection}

\noindent The convergence metric $\sigma^{(r)}$ measures the average pairwise similarity between agent messages at round $r$:

$$\sigma^{(r)} = \frac{2}{n(n-1)} \sum_{i<j} \mathrm{sim}\bigl(\mathbf{e}(m_i^{(r)}),\, \mathbf{e}(m_j^{(r)})\bigr)$$

\noindent Reasoning terminates at round $r^*$ when any of:
\begin{itemize}[nosep,leftmargin=2em]
\item $\sigma^{(r^*)} \geq \tau_{\mathrm{conv}}$ (agents have reached consensus);
\item $|\sigma^{(r^*)} - \sigma^{(r^*-1)}| < \epsilon_{\mathrm{delta}}$ (convergence has stalled);
\item $r^* = R_{\max}$ (maximum rounds reached);
\item $C_{\mathrm{cum}} \geq C_{\max}$ (cost budget exceeded).
\end{itemize}

\subsection*{5.5 MasterObserver Architecture}

\noindent The MasterObserver $\mathcal{O}$ implements a \textbf{read-only event bus pattern}:

\begin{enumerate}[nosep,leftmargin=2em]
\item All messages between agents are copied to $\mathcal{O}$ without blocking the sender or receiver;
\item $\mathcal{O}$ records each message as a tuple $(a_i, a_j, m, r, t)$ in the trace $\mathcal{T}$;
\item After each round, $\mathcal{O}$ computes:
  \begin{itemize}[nosep,leftmargin=1.5em]
  \item Health score $h^{(r)}$ based on message count, average length, and error rate;
  \item Per-agent contribution score $\kappa_i$;
  \item Anomaly flags for silent agents ($|m_i| = 0$) or circular reasoning;
  \end{itemize}
\item $\mathcal{O}$ has \textbf{zero inference cost} --- it observes but never calls a language model.
\end{enumerate}

\subsection*{5.6 Backend Fallback and Cost Control}

\noindent The backend fallback chain supports three categories:

\begin{center}
\begin{tabular}{lll}
\toprule
\textbf{Category} & \textbf{Examples} & \textbf{Cost Rate} \\
\midrule
Cloud API & Anthropic Claude, OpenAI GPT & \$0.25--15/MTok \\
Local GPU & Ollama (Qwen2.5, Llama), vLLM & \$0.0001/MTok \\
Adapter & LoRA-augmented local models & \$0.0001/MTok \\
\bottomrule
\end{tabular}
\end{center}

\noindent Cost budget enforcement tracks $C_{\mathrm{cum}}$ across all inference calls. When $C_{\mathrm{cum}} \geq C_{\max}$, reasoning halts gracefully, returning partial results with a \texttt{budget\_exceeded} flag.

\subsection*{5.7 Hierarchical Aggregation}

\noindent Final-round messages are aggregated using the graph topology:

\begin{enumerate}[nosep,leftmargin=2em]
\item Compute betweenness centrality $\mathrm{BC}(v_i)$ for each node in $G^*$;
\item Select the highest-centrality node as root;
\item Construct a BFS spanning tree from root;
\item Aggregate bottom-up: leaf responses $\to$ hub syntheses $\to$ root answer.
\end{enumerate}

\noindent This preserves the knowledge graph's structural hierarchy in the final answer, unlike flat concatenation approaches.

\vspace{8mm}

%% =============================================================================
%% SECTION 6: INDUSTRIAL APPLICABILITY
%% =============================================================================

\section*{Industrial Applicability}

\noindent The CogniGraph system is applicable to any domain requiring structured reasoning over large knowledge graphs, including:

\begin{itemize}[nosep,leftmargin=2em]
\item \textbf{Regulatory compliance}: Reasoning over legal knowledge graphs (EU AI Act, financial regulations) where multi-hop reasoning across articles and cross-references is required;
\item \textbf{Scientific research}: Reasoning over biomedical knowledge graphs, materials science ontologies, or academic citation networks;
\item \textbf{Enterprise knowledge management}: Multi-agent reasoning over corporate knowledge bases with role-based access and audit trails;
\item \textbf{Multi-hop question answering}: Answering complex questions that require synthesizing information from multiple knowledge sources.
\end{itemize}

\vspace{2mm}

\noindent The system is implemented as an open-source Python SDK (\texttt{pip install cognigraph}) with CLI (\texttt{kogni}), REST API server, and connectors for NetworkX, Neo4j, and JSON graph formats.

\vspace{15mm}

\noindent\rule{\textwidth}{0.4pt}

\begin{center}
{\small\itshape Quantamix Solutions B.V., Buitendijks\,2, 1422\,MM Uithoorn, The Netherlands\\
European Patent Application --- Continuation-in-Part of EP26162901.8\\
Reference: COGNIGRAPH-CIP-NL-2026-001}
\end{center}

\end{document}
